% ============================================================
%  GUÍA 2: Plano Cartesiano, Relaciones y Gráficas Básicas
%  Curso de Introducción a la Universidad — Precálculo y Cálculo I
%  Autor: Ing. Gabriel Astudillo
%  Clase cubierta: 2
% ============================================================

\documentclass[12pt, a4paper]{article}

% ── Paquetes ─────────────────────────────────────────────────
\usepackage{mathwizards-guia}
\mwguia{Guía 2 — P. Cartesiano, Relaciones y Funciones}{Ing. Gabriel Astudillo $\cdot$ Curso Preuniversitario}

\begin{document}

% ── Portada / Título ─────────────────────────────────────────
\mwportada{Guía de Estudio 2}{Plano Cartesiano, Relaciones y Gráficas Básicas}{Distancia, Punto Medio, Relaciones, Funciones, Dominio, Rango y Gráficas Elementales}

\vspace{4pt}
\begin{cuadroinfo}
  \small
  \textbf{Presentación:} Esta guía introduce el lenguaje geométrico y analítico de las funciones. Aprenderás a calcular distancias y puntos medios, a reconocer cuándo una relación es función, a determinar dominios y rangos con rigor y a trazar e interpretar las gráficas elementales: lineales, cuadráticas, cúbicas, de valor absoluto y raíces. Cada sección incluye teoría, ejercicios resueltos paso a paso y ejercicios propuestos con su clave.
  \\[4pt]
  \textbf{Nivel de Dificultad:}\quad
  \facil{} Básico / Directo \quad
  \medio{} Intermedio \quad
  \dificil{} Desafiante \quad
  \muyDificil{} Muy Difícil / Reto
\end{cuadroinfo}

\vspace{8pt}

\section{Plano Cartesiano: Distancia y Punto Medio}
% ============================================================

\subsection{Fórmulas fundamentales}

\begin{cuadroteoria}
  Dados dos puntos $P_1(x_1, y_1)$ y $P_2(x_2, y_2)$ en el plano cartesiano:

  \textbf{Distancia entre dos puntos:}
  $$d(P_1, P_2) = \sqrt{(x_2 - x_1)^2 + (y_2 - y_1)^2}$$

  \textbf{Punto medio del segmento:}
  $$M = \left(\frac{x_1 + x_2}{2}, \; \frac{y_1 + y_2}{2}\right)$$
\end{cuadroteoria}

\newpage %gabo

\subsection{Ejercicios resueltos}

\textbf{Ejemplo R1.} Calcula la distancia entre $A(-2, 3)$ y $B(4, -5)$, y el punto medio del segmento $AB$.

\begin{cuadrosolucion}
  \textbf{Solución:}
  \begin{itemize}
    \item Distancia: $d = \sqrt{(4 - (-2))^2 + (-5 - 3)^2} = \sqrt{6^2 + (-8)^2} = \sqrt{36 + 64} = \sqrt{100} = 10$.
    \item Punto medio: $M = \left(\dfrac{-2 + 4}{2}, \; \dfrac{3 + (-5)}{2}\right) = (1, -1)$.
  \end{itemize}
\end{cuadrosolucion}

\textbf{Ejemplo R2.} Los puntos $A(1, 2)$, $B(4, 6)$ y $C(7, 2)$ forman un triángulo. Determina si es isósceles, equilátero o escaleno.

\begin{cuadrosolucion}
  \textbf{Solución:} Calculamos las tres distancias:
  \begin{itemize}
    \item $AB = \sqrt{(4-1)^2 + (6-2)^2} = \sqrt{9 + 16} = 5$.
    \item $BC = \sqrt{(7-4)^2 + (2-6)^2} = \sqrt{9 + 16} = 5$.
    \item $AC = \sqrt{(7-1)^2 + (2-2)^2} = \sqrt{36} = 6$.
  \end{itemize}
  Como $AB = BC = 5$ y $AC = 6$, el triángulo es \textbf{isósceles}.
\end{cuadrosolucion}

\textbf{Ejemplo R3.} Uno de los extremos de un segmento es $A(2, -1)$ y su punto medio es $M(5, 3)$. Encuentra el otro extremo $B$.

\begin{cuadrosolucion}
  \textbf{Solución:} Usamos la fórmula del punto medio, donde $M = \left(\dfrac{x_A + x_B}{2}, \dfrac{y_A + y_B}{2}\right)$:
  \begin{itemize}
    \item $5 = \dfrac{2 + x_B}{2} \Rightarrow 10 = 2 + x_B \Rightarrow x_B = 8$.
    \item $3 = \dfrac{-1 + y_B}{2} \Rightarrow 6 = -1 + y_B \Rightarrow y_B = 7$.
  \end{itemize}
  Por tanto, $B(8, 7)$.
\end{cuadrosolucion}

\subsection{Ejercicios propuestos}

\begin{enumerate}[series=ejercicios, leftmargin=*]
  \item \facil{} Calcula la distancia entre $P(3, 4)$ y $Q(-2, 1)$.

  \item \facil{} Calcula el punto medio del segmento que une $A(-6, 8)$ y $B(4, -2)$.

  \item \facil{} Determina si los puntos $A(0, 0)$, $B(6, 8)$ y $C(12, 0)$ forman un triángulo isósceles.

  \item \medio{} Los puntos $A(-2, 1)$, $B(1, 5)$ y $C(6, -2)$ forman un triángulo. Calcula su perímetro.

  \item \medio{} Encuentra todos los puntos sobre el eje $x$ que están a una distancia de $5$ del punto $P(3, 4)$.

  \item \medio{} El punto medio del segmento $AB$ es $M(-1, 4)$. Si $A = (3, 2)$, encuentra $B$.

  \item \dificil{} Demuestra que los puntos $A(0, 3)$, $B(4, 0)$, $C(8, 3)$ y $D(4, 6)$ forman un rombo (todos sus lados son iguales).

  \item \dificil{} Un triángulo tiene vértices $A(2, 3)$, $B(6, 7)$ y $C(10, 3)$. Calcula la longitud de la mediana que va desde $B$ hasta el punto medio de $AC$.
\end{enumerate}

% ============================================================

\section{Relaciones y Funciones}
% ============================================================

\subsection{Conceptos fundamentales}

\begin{cuadroteoria}
  \textbf{Relación:} Un conjunto de pares ordenados $(x, y)$. Ejemplo: $R = \{(1, 2), (1, 3), (2, 5)\}$.

  \textbf{Función:} Una relación donde cada $x$ (del dominio) se asocia con \textbf{exactamente un} $y$ (del rango).\\[6pt]

  En otras palabras, para una función no pueden existir dos pares con el mismo $x$ y distintos $y$. Se escribe $y = f(x)$.
\end{cuadroteoria}

\begin{cuadrosolucion}
  \textbf{Prueba de la recta vertical:} Una gráfica representa una función si \textbf{toda recta vertical} la intersecta a lo sumo una vez.

  \textbf{Dominio:} Conjunto de todos los valores de entrada $x$ para los cuales la función está definida.

  \textbf{Rango (o imagen):} Conjunto de todos los valores de salida $y = f(x)$ que la función produce.
\end{cuadrosolucion}

\begin{cuadroadvertencia}
  \textbf{Reglas para determinar el dominio:}
  \begin{itemize}[leftmargin=*]
    \item Denominadores: el denominador nunca puede ser cero.
    \item Raíces pares ($\sqrt{\phantom{x}}$, $\sqrt[4]{\phantom{x}}$, etc.): el radicando debe ser $\geq 0$.
    \item Raíces impares ($\sqrt[3]{\phantom{x}}$, $\sqrt[5]{\phantom{x}}$, etc.): sin restricción.
    \item Logaritmos: el argumento debe ser $> 0$ (aunque esto se verá en la Guía 4).
  \end{itemize}
\end{cuadroadvertencia}

\subsection{Ejercicios resueltos}

\textbf{Ejemplo R1.} Determina cuáles de las siguientes relaciones son funciones:
\begin{enumerate}[label=\alph*)]
  \item $R_1 = \{(1, 2), (2, 3), (3, 4), (4, 5)\}$
  \item $R_2 = \{(1, 2), (1, 3), (2, 4)\}$
  \item $R_3 = \{(x, y) : y = x^2\}$
  \item $R_4 = \{(x, y) : x = y^2\}$
\end{enumerate}

\newpage %gabo

\begin{cuadrosolucion}
  \textbf{Solución:}
  \begin{itemize}
    \item $R_1$ es función: cada $x$ tiene un único $y$.
    \item $R_2$ \textbf{no} es función: el valor $x = 1$ se asocia con dos $y$ distintos ($2$ y $3$).
    \item $R_3$ es función: para cada $x$, hay un único $x^2$.
    \item $R_4$ \textbf{no} es función: para $x = 4$ hay dos $y$ posibles, $y = 2$ y $y = -2$.
  \end{itemize}
\end{cuadrosolucion}

\textbf{Ejemplo R2.} Encuentra el dominio de $f(x) = \dfrac{x + 2}{x^2 - 9}$.

\begin{cuadrosolucion}
  \textbf{Solución:} El denominador debe ser distinto de cero:
  $$x^2 - 9 \neq 0 \Rightarrow (x - 3)(x + 3) \neq 0 \Rightarrow x \neq 3,\; x \neq -3.$$
  Dominio: $(-\infty, -3) \cup (-3, 3) \cup (3, \infty)$.
\end{cuadrosolucion}

\textbf{Ejemplo R3.} Encuentra el dominio de $f(x) = \sqrt{2x - 6}$.

\begin{cuadrosolucion}
  \textbf{Solución:} El radicando debe ser no negativo:
  $$2x - 6 \geq 0 \Rightarrow 2x \geq 6 \Rightarrow x \geq 3.$$
  Dominio: $[3, \infty)$.
\end{cuadrosolucion}

\textbf{Ejemplo R4.} Encuentra el dominio y el rango de $f(x) = x^2 - 4$.

\begin{cuadrosolucion}
  \textbf{Solución:}
  \begin{itemize}
    \item \textbf{Dominio:} No hay restricciones (no hay denominador ni raíz par). Dominio: $(-\infty, \infty)$.
    \item \textbf{Rango:} Como $x^2 \geq 0$ para todo $x$, entonces $x^2 - 4 \geq -4$. El valor mínimo se alcanza en $x = 0$. Rango: $[-4, \infty)$.
  \end{itemize}
\end{cuadrosolucion}

\textbf{Ejemplo R5.} Encuentra el dominio de $f(x) = \dfrac{1}{\sqrt{x - 1}}$.

\begin{cuadrosolucion}
  \textbf{Solución:} Tenemos dos restricciones:
  \begin{itemize}
    \item El radicando debe ser $\geq 0$: $x - 1 \geq 0$, es decir, $x \geq 1$.
    \item El denominador no puede ser cero: $\sqrt{x - 1} \neq 0$, es decir, $x - 1 \neq 0$, o sea $x \neq 1$.
  \end{itemize}
  Combinando: Dominio $= (1, \infty)$.
\end{cuadrosolucion}

\newpage

\subsection{Ejercicios propuestos}

\begin{enumerate}[resume=ejercicios, leftmargin=*]
  \item \facil{} Determina cuáles de las siguientes son funciones (justifica):
        \begin{enumerate}[label=\alph*)]
          \item $\{(2, 3), (3, 4), (4, 5)\}$
          \item $\{(2, 3), (2, 4), (5, 6)\}$
          \item $y = 2x + 1$
          \item $x^2 + y^2 = 25$ (circunferencia)
        \end{enumerate}

  \item \facil{} Encuentra el dominio: $f(x) = \dfrac{3}{x - 5}$

  \item \facil{} Encuentra el dominio: $f(x) = \sqrt{x + 7}$

  \item \facil{} Encuentra el dominio: $f(x) = x^3 - 2x + 1$

  \item \medio{} Encuentra el dominio: $f(x) = \dfrac{x}{x^2 - 4}$

  \item \medio{} Encuentra el dominio: $f(x) = \sqrt{9 - x^2}$

  \item \medio{} Encuentra el dominio: $f(x) = \dfrac{2x + 1}{\sqrt{x + 3}}$

  \item \medio{} Encuentra el dominio y el rango de $f(x) = |x| - 2$.

  \item \medio{} Encuentra el dominio y el rango de $f(x) = \sqrt{x - 1} + 3$.

  \item \dificil{} Encuentra el dominio de $f(x) = \dfrac{\sqrt{x + 2}}{x - 1}$.

  \item \dificil{} Encuentra el dominio de $f(x) = \sqrt{\dfrac{x - 2}{x + 3}}$.

  \item \dificil{} El costo de producir $x$ unidades de un producto es $C(x) = 2x^2 + 50x + 100$ dólares. ¿Cuál es el dominio realista de esta función en el contexto del problema? ¿Cuál es el costo de producir 10 unidades?
\end{enumerate}

\newpage

% ============================================================

\section{Gráficas de Funciones Básicas}
% ============================================================

\subsection{Las funciones elementales}

\begin{cuadroteoria}
  \textbf{Funciones básicas que debes dominar:}

  \begin{center}
    \begin{tabular}{l c p{7.2cm}}
      \toprule
      \textbf{Función} & \textbf{Ecuación}      & \textbf{Características}                          \\
      \midrule
      Lineal           & $f(x) = mx + b$        & Línea recta; pendiente $m$; intercepto $(0, b)$   \\
      Cuadrática       & $f(x) = ax^2 + bx + c$ & Parábola; vértice en $x = -\frac{b}{2a}$          \\
      Cúbica           & $f(x) = x^3$           & Pasa por el origen; punto de inflexión en $(0,0)$ \\
      Valor absoluto   & $f(x) = |x|$           & Forma de V; vértice en el origen                  \\
      Raíz cuadrada    & $f(x) = \sqrt{x}$      & Solo para $x \geq 0$; creciente                   \\
      Recíproca        & $f(x) = \dfrac{1}{x}$  & Hipérbola; asíntotas en $x = 0$ e $y = 0$         \\
      \bottomrule
    \end{tabular}
  \end{center}
\end{cuadroteoria}

\begin{cuadrosolucion}
  \textbf{Interceptos:}
  \begin{itemize}[leftmargin=*]
    \item \textbf{Intercepto $y$:} se encuentra evaluando $f(0)$. Es el punto $(0, f(0))$.
    \item \textbf{Interceptos $x$ (raíces o ceros):} se encuentran resolviendo $f(x) = 0$. Son los puntos $(x, 0)$.
  \end{itemize}
\end{cuadrosolucion}

\begin{cuadroadvertencia}
  \textbf{Vértice de una parábola:} Para $f(x) = ax^2 + bx + c$:
  $$x_v = -\frac{b}{2a}, \qquad y_v = f(x_v)$$
  \begin{itemize}
    \item Si $a > 0$: la parábola abre hacia arriba y el vértice es un mínimo.
    \item Si $a < 0$: la parábola abre hacia abajo y el vértice es un máximo.
  \end{itemize}
\end{cuadroadvertencia}

\newpage

\subsection{Ejercicios resueltos}

\textbf{Ejemplo R1.} Encuentra los interceptos y el vértice de $f(x) = x^2 - 4x + 3$, y esboza la gráfica.

\begin{cuadrosolucion}
  \textbf{Solución:}
  \begin{itemize}
    \item \textbf{Intercepto $y$:} $f(0) = 0 - 0 + 3 = 3$. Punto $(0, 3)$.
    \item \textbf{Interceptos $x$:} $x^2 - 4x + 3 = 0 \Rightarrow (x - 1)(x - 3) = 0 \Rightarrow x = 1$ y $x = 3$. Puntos $(1, 0)$ y $(3, 0)$.
    \item \textbf{Vértice:} $x_v = -\dfrac{-4}{2(1)} = 2$; $y_v = f(2) = 4 - 8 + 3 = -1$. Vértice: $(2, -1)$.
  \end{itemize}
  \begin{center}
    \begin{tikzpicture}[scale=0.7]
      \draw[thin, gray!50] (-1,-3) grid (5,5);
      \draw[thick, -Latex] (-1.5,0) -- (5.5,0) node[right] {$x$};
      \draw[thick, -Latex] (0,-3.5) -- (0,5.5) node[above] {$y$};
      \draw[domain=-0.5:4.5, smooth, very thick, mwprimario] plot (\x, {\x*\x - 4*\x + 3});
      \filldraw[red] (0,3) circle (2.5pt) node[above right] {$(0,3)$};
      \filldraw[red] (1,0) circle (2.5pt) node[below] {$(1,0)$};
      \filldraw[red] (3,0) circle (2.5pt) node[below right] {$(3,0)$};
      \filldraw[red] (2,-1) circle (2.5pt) node[below] {$(2,-1)$};
      \draw[gray] (0,0) node[below left] {$0$};
    \end{tikzpicture}
  \end{center}
\end{cuadrosolucion}

\textbf{Ejemplo R2.} Determina el vértice de $f(x) = -x^2 + 6x - 5$ y di si es máximo o mínimo.

\begin{cuadrosolucion}
  \textbf{Solución:} $a = -1$, $b = 6$, $c = -5$.
  $$x_v = -\frac{6}{2(-1)} = 3; \qquad y_v = f(3) = -(3)^2 + 6(3) - 5 = -9 + 18 - 5 = 4.$$
  Vértice: $(3, 4)$. Como $a < 0$, el vértice es un \textbf{máximo}.
\end{cuadrosolucion}

\textbf{Ejemplo R3.} Encuentra los interceptos de $f(x) = |x + 2| - 1$.

\begin{cuadrosolucion}
  \textbf{Solución:}
  \begin{itemize}
    \item \textbf{Intercepto $y$:} $f(0) = |0 + 2| - 1 = 2 - 1 = 1$. Punto $(0, 1)$.
    \item \textbf{Interceptos $x$:} $|x + 2| - 1 = 0 \Rightarrow |x + 2| = 1$.
          \begin{itemize}
            \item $x + 2 = 1 \Rightarrow x = -1$.
            \item $x + 2 = -1 \Rightarrow x = -3$.
          \end{itemize}
          Puntos $(-1, 0)$ y $(-3, 0)$.
  \end{itemize}
\end{cuadrosolucion}

\newpage

\subsection{Ejercicios propuestos}

\begin{enumerate}[resume=ejercicios, leftmargin=*]
  \item \facil{} Encuentra los interceptos y el vértice de $f(x) = x^2 - 2x - 8$.

  \item \facil{} Encuentra los interceptos de $f(x) = 2x - 6$.

  \item \facil{} Determina el vértice de $f(x) = x^2 + 6x + 5$ y clasifícalo.

  \item \facil{} Encuentra los interceptos de $f(x) = |x - 3| - 2$.

  \item \medio{} Determina el vértice de $f(x) = -2x^2 + 8x - 3$ y los interceptos.

  \item \medio{} Encuentra los interceptos de $f(x) = 9 - x^2$.

  \item \medio{} Encuentra los interceptos de $f(x) = x^3 - 4x$.

  \item \medio{} Dada $f(x) = x^2 + bx + 4$, encuentra el valor de $b$ sabiendo que el vértice está en $x = 3$.

  \item \dificil{} Una parábola tiene vértice en $(2, 5)$ y pasa por el punto $(4, -3)$. Encuentra su ecuación de la forma $f(x) = a(x - h)^2 + k$.

  \item \dificil{} Encuentra los valores de $k$ para que la parábola $f(x) = kx^2 + 2x + 1$ corte el eje $x$ en dos puntos distintos.

  \item \dificil{} Un rectángulo tiene su base sobre el eje $x$ y sus vértices superiores sobre la parábola $y = 12 - x^2$. Encuentra el área máxima posible del rectángulo.

  \item \dificil{} Determina el conjunto de puntos $(x, y)$ que cumplen $y = |x| - x$ y esboza la gráfica. (Pista: considera casos $x \geq 0$ y $x < 0$.)
\end{enumerate}

% ============================================================

% ──────────────────────────────────────────────────────────
%  NUEVOS EJERCICIOS PROPUESTOS (18) — INSERCIÓN MANUAL
% ──────────────────────────────────────────────────────────
\subsection{Ejercicios propuestos: Repaso de Distancia, Relaciones y Dominios}

\begin{enumerate}[resume=ejercicios, leftmargin=*]
  \item \facil{} Calcula la distancia entre los puntos $A(-3, 5)$ y $B(4, -2)$.

  \item \facil{} Halla el punto medio del segmento cuyos extremos son $M(-6, 8)$ y $N(2, -4)$.

  \item \facil{} Determina el dominio de la función:
        $$f(x) = \sqrt{4 - x^2}$$

  \item \medio{} Determina cuáles de las siguientes curvas representan una función (justifica con la prueba de la recta vertical):
        $$a)\; y = x^2 \qquad\qquad\qquad b)\; x = y^2$$

  \item \medio{} Halla el dominio y el rango de la función racional (simplifica antes):
        $$f(x) = \frac{x + 1}{x^2 - 1}$$

  \item \medio{} Calcula la distancia del punto $P(2, -3)$ al origen y al punto $Q(-1, 5)$.

  \item \medio{} Halla el punto medio del segmento con extremos $A(3a, -2b)$ y $B(-a, 4b)$.

  \item \dificil{} Determina el dominio de la función:
        $$f(x) = \sqrt{x^2 - 5x + 6}$$

  \item \dificil{} Halla los puntos del eje $y$ que se encuentran a distancia $5$ del punto $(3, 0)$.

  \item \dificil{} Determina el dominio y el rango de:
        $$f(x) = \frac{x^2 - 1}{x^2 + 1}$$
\end{enumerate}

\subsection{Ejercicios propuestos: Gráficas Elementales Avanzadas}

\begin{enumerate}[resume=ejercicios, leftmargin=*]
  \item \muyDificil{} Determina el dominio de la función:
        $$f(x) = \frac{\sqrt{x + 2}}{|x| - 1}$$

  \item \facil{} Halla los interceptos con los ejes de la parábola $y = x^2 - 4x + 3$.

  \item \medio{} Halla el vértice de la parábola $y = -2x^2 + 8x - 5$.

  \item \medio{} Determina la ecuación de la recta que pasa por los puntos $(1, 2)$ y $(3, 6)$.

  \item \medio{} Describe la transformación que produce la gráfica de $y = |x - 2| + 1$ a partir de $y = |x|$ e indica su vértice.

  \item \dificil{} ¿La circunferencia $x^2 + y^2 = 9$ representa una función? Justifica tu respuesta.

  \item \dificil{} Determina el dominio y el rango de $f(x) = \sqrt{9 - x^2}$ e identifica su gráfica.

  \item \muyDificil{} Halla la ecuación de la función cuadrática cuyo vértice es $(2, -1)$ y que pasa por el punto $(0, 3)$.
\end{enumerate}
\newpage
% ============================================================
\ifmwclaves
\section{Clave de Respuestas}
% ============================================================

\subsection*{Sección 1: Distancia y Punto Medio (Ejercicios 1--8)}

\begin{enumerate}[series=respuestas, leftmargin=*, label=\textbf{\arabic*.}]
  \item $d = \sqrt{(-2-3)^2 + (1-4)^2} = \sqrt{25 + 9} = \sqrt{34}$

  \item $M = \left(\dfrac{-6+4}{2}, \dfrac{8+(-2)}{2}\right) = (-1, 3)$

  \item $AB = \sqrt{36 + 64} = 10$; $BC = \sqrt{36 + 64} = 10$; $AC = 12$. Sí, es isósceles ($AB = BC$).

  \item $AB = \sqrt{9 + 16} = 5$; $BC = \sqrt{25 + 49} = \sqrt{74}$; $AC = \sqrt{64 + 9} = \sqrt{73}$. Perímetro: $5 + \sqrt{74} + \sqrt{73}$.

  \item Sea $Q(x, 0)$ en el eje $x$: $d(Q, P) = \sqrt{(x - 3)^2 + 16} = 5$. Elevando al cuadrado: $(x-3)^2 = 9$, $x - 3 = \pm 3$, $x = 0$ o $x = 6$. Puntos: $(0, 0)$ y $(6, 0)$.

  \item $B = (2(-1) - 3, \; 2(4) - 2) = (-5, 6)$.

  \item $AB = \sqrt{16 + 9} = 5$; $BC = \sqrt{16 + 9} = 5$; $CD = \sqrt{16 + 9} = 5$; $DA = \sqrt{16 + 9} = 5$. Todos los lados miden 5, es un rombo.

  \item Punto medio de $AC$: $M = \left(\dfrac{2 + 10}{2}, \dfrac{3 + 3}{2}\right) = (6, 3)$. Mediana desde $B(6, 7)$ a $M(6, 3)$: $d = \sqrt{0 + 16} = 4$.
\end{enumerate}

\newpage


\subsection*{Sección 2: Relaciones y Funciones (Ejercicios 9--20)}

\begin{enumerate}[resume=respuestas, leftmargin=*, label=\textbf{\arabic*.}]
  \item a) Función. b) No es función (ejemplo: $x=2$ con dos $y$). c) Función (recta). d) No es función (para un $x$ hay dos $y$; falla la recta vertical).

  \item $(-\infty, 5) \cup (5, \infty)$

  \item $[-7, \infty)$

  \item $(-\infty, \infty)$

  \item $(-\infty, -2) \cup (-2, 2) \cup (2, \infty)$

  \item $9 - x^2 \geq 0 \Rightarrow x^2 \leq 9 \Rightarrow [-3, 3]$

  \item $x + 3 > 0 \Rightarrow x > -3$ (además el denominador $\sqrt{x+3} \neq 0$ ya está cubierto). Dominio: $(-3, \infty)$.

  \item Dominio: $(-\infty, \infty)$; rango: $[-2, \infty)$.

  \item Dominio: $[1, \infty)$; rango: $[3, \infty)$.

  \item $x + 2 \geq 0 \Rightarrow x \geq -2$, y $x - 1 \neq 0 \Rightarrow x \neq 1$. Dominio: $[-2, 1) \cup (1, \infty)$.

  \item Resolvemos la desigualdad $\dfrac{x-2}{x+3} \geq 0$: puntos clave $x = -3$ y $x = 2$. Dominio: $(-\infty, -3) \cup [2, \infty)$ (excluimos $x = -3$ porque el denominador se anula).

  \item Dominio realista: $x \geq 0$ y $x$ entero (no se pueden producir cantidades negativas o fraccionarias). $C(10) = 2(100) + 50(10) + 100 = 200 + 500 + 100 = 800$ dólares.
\end{enumerate}

\newpage


\subsection*{Sección 3: Gráficas de Funciones Básicas (Ejercicios 21--32)}

\begin{enumerate}[resume=respuestas, leftmargin=*, label=\textbf{\arabic*.}]
  \item Intercepto $y$: $(0, -8)$. Interceptos $x$: $x^2 - 2x - 8 = (x - 4)(x + 2) = 0 \Rightarrow x = 4$, $x = -2$. Vértice: $x_v = 1$, $y_v = -9$; vértice $(1, -9)$.

  \item Intercepto $y$: $(0, -6)$. Intercepto $x$: $2x - 6 = 0 \Rightarrow x = 3$; punto $(3, 0)$.

  \item $x_v = -3$, $y_v = f(-3) = 9 - 18 + 5 = -4$. Vértice: $(-3, -4)$; como $a = 1 > 0$, es un mínimo.

  \item Intercepto $y$: $f(0) = 1$, punto $(0, 1)$. Interceptos $x$: $|x - 3| = 2 \Rightarrow x = 5$ o $x = 1$; puntos $(5, 0)$ y $(1, 0)$.

  \item Vértice: $x_v = 2$, $y_v = f(2) = -8 + 16 - 3 = 5$; vértice $(2, 5)$ máximo. Intercepto $y$: $(0, -3)$. Interceptos $x$: $x = \dfrac{-8 \pm \sqrt{64 - 24}}{-4} = \dfrac{-8 \pm \sqrt{40}}{-4} = \dfrac{-8 \pm 2\sqrt{10}}{-4} = \dfrac{4 \mp \sqrt{10}}{2}$. Aproximadamente $x \approx 0{,}42$ y $x \approx 3{,}58$.

  \item Intercepto $y$: $(0, 9)$. Interceptos $x$: $9 - x^2 = 0 \Rightarrow x = \pm 3$; puntos $(3, 0)$ y $(-3, 0)$.

  \item Intercepto $y$: $(0, 0)$. Interceptos $x$: $x^3 - 4x = x(x^2 - 4) = x(x-2)(x+2) = 0$; $x = 0$, $x = 2$, $x = -2$.

  \item Vértice en $x_v = -\frac{b}{2} = 3 \Rightarrow b = -6$.

  \item Vértice $(2, 5)$: $f(x) = a(x - 2)^2 + 5$. Como pasa por $(4, -3)$: $-3 = a(4-2)^2 + 5 \Rightarrow -3 = 4a + 5 \Rightarrow 4a = -8 \Rightarrow a = -2$. Ecuación: $f(x) = -2(x - 2)^2 + 5$.

  \item Discriminante $D = 4 - 4k(1) = 4 - 4k > 0 \Rightarrow k < 1$ (con $k \neq 0$, de lo contrario no es cuadrática).

  \item El rectángulo tiene base $2x$ y altura $12 - x^2$ (con $0 < x < \sqrt{12}$). Área: $A(x) = 2x(12 - x^2) = 24x - 2x^3$. Para maximizar con cálculo próximo, derivamos (esto se verá en la Guía 7): $A'(x) = 24 - 6x^2 = 0 \Rightarrow x = 2$ (positivo). $A_{\max} = 24(2) - 2(8) = 48 - 16 = 32$ unidades cuadradas.

  \item Para $x \geq 0$: $y = x - x = 0$ (el semieje $x$ positivo). Para $x < 0$: $y = -x - x = -2x$ (recta con pendiente $-2$ en la región negativa). La gráfica es: el semieje positivo $x$ más la semi-recta $y = -2x$ para $x < 0$.
\end{enumerate}

\newpage


% ──────────────────────────────────────────────────────────
%  NUEVAS RESPUESTAS (18) — INSERCIÓN MANUAL
% ──────────────────────────────────────────────────────────
\subsection*{Sección 4: Repaso de Distancia, Relaciones y Dominios (Ejercicios 33--42)}
\begin{enumerate}[resume=respuestas, leftmargin=*, label=\textbf{\arabic*.}]
  \item $d = \sqrt{(4 + 3)^2 + (-2 - 5)^2} = \sqrt{49 + 49} = 7\sqrt{2}$
  \item $M = \left(\dfrac{-6 + 2}{2}, \dfrac{8 - 4}{2}\right) = (-2, 2)$
  \item $4 - x^2 \ge 0 \Rightarrow x^2 \le 4 \Rightarrow$ dominio $[-2, 2]$
  \item a) Sí es función (toda recta vertical corta una sola vez). \quad b) No es función (para un $x$ hay dos $y$; falla la recta vertical).
  \item Simplificando: $f(x) = \dfrac{1}{x - 1}$ (con $x \neq -1$). Dominio: $\mathbb{R} \setminus \{\pm 1\}$; \quad rango: $\mathbb{R} \setminus \{0, -1/2\}$.
  \item $d(P, O) = \sqrt{13}$ \quad y \quad $d(P, Q) = \sqrt{73}$
  \item $M = (a, b)$
  \item $x^2 - 5x + 6 \ge 0 \Rightarrow (x - 2)(x - 3) \ge 0 \Rightarrow$ dominio $(-\infty, 2] \cup [3, \infty)$
  \item $(0, 4)$ y $(0, -4)$ \quad (de $9 + y^2 = 25$)
  \item Dominio: $\mathbb{R}$. \quad Rango: $[-1, 1)$ (despejando $x^2 = \dfrac{1 + y}{1 - y} \ge 0$)
\end{enumerate}

\subsection*{Sección 5: Gráficas Elementales Avanzadas (Ejercicios 43--50)}
\begin{enumerate}[resume=respuestas, leftmargin=*, label=\textbf{\arabic*.}]
  \item Dominio: $[-2, -1) \cup (-1, 1) \cup (1, \infty)$ \quad ($x \ge -2$ y $x \neq \pm 1$)
  \item Interceptos con $x$: $(1, 0)$ y $(3, 0)$; \quad intercepto con $y$: $(0, 3)$
  \item $x_v = 2$, $y_v = 3$; \quad vértice $(2, 3)$ (máximo)
  \item $y = 2x$
  \item Traslación 2 unidades a la derecha y 1 hacia arriba; vértice $(2, 1)$; forma de $V$ hacia arriba.
  \item No es función: toda recta vertical $x = c$ con $-3 < c < 3$ corta a la circunferencia en dos puntos.
  \item Dominio $[-3, 3]$; rango $[0, 3]$; gráfica: semicircunferencia superior de radio 3.
  \item $y = a(x - 2)^2 - 1$; con $(0, 3)$: $3 = 4a - 1 \Rightarrow a = 1$; \quad $y = x^2 - 4x + 3$
\end{enumerate}
\fi
\end{document}
